\section{Related work}
\label{sec:related}

\paragraph{Video token reduction.} Run-length tokenization \citep{choudhury2024rlt} compares each spatio-temporal patch with the one at the same position in the previous frame by an $L^1$ pixel difference and merges runs of similar patches into one token with a length encoding. Efficient Video Sampling \citep{evs2025} and TimeChat-Online \citep{yao2025timechatonline} threshold a per-patch difference between consecutive frames inside a video LLM. DyCoke \citep{tao2025dycoke} merges tokens at the same position across groups of four sampled frames by cosine similarity of their encoder features, keeping the first frame of each group, and prunes further by attention in the LLM. FrameFusion \citep{fu2024framefusion} merges adjacent-frame tokens by cosine similarity in early LLM layers and prunes by importance after. PruneVid \citep{huang2025prunevid} segments the video into scenes and merges the tokens that are static across a segment, then clusters spatially. LongVU \citep{shen2024longvu} drops frames whose DINOv2 features are similar to the previous frame's and pools the rest spatially. FastV \citep{chen2024fastv}, VisionZip \citep{yang2025visionzip} and SparseVLM \citep{zhang2024sparsevlm} prune by attention within an image and have no temporal criterion. Section~\ref{sec:audit} re-implements the temporal decision rule of each of the first six in pixel space and says exactly what is kept of the published method; the attention-based methods need the model and are out of scope here. Trajectory tokenization \citep{zhang2025trajectory} assigns one token per panoptic sub-object trajectory, which is the entity-token idea without a certificate; the moving-camera certificate of Section~\ref{sec:theory} is what would let such a token stay frozen with a proof.

\paragraph{Conditional replenishment and deadband sampling.} Sending a sample only when it has drifted by more than a threshold from the last sample sent is conditional replenishment \citep{mounts1969} in video coding and send-on-delta sampling \citep{miskowicz2006} in event-based control. The anchored rule is this idea with a sup-norm statistic; motion-compensated prediction with a residual test is every codec since H.261. What is new in \citet{nair2026certifiedskip} and here is the certificate, the quotient by a motion group, and the competitive analysis of Section~\ref{sec:opt}, which to our knowledge has not been stated for deadband schemes.

\paragraph{Competitive analysis.} Theorem~\ref{thm:opt} is an online-algorithms statement in the sense of \citet{borodin1998}: a causal rule against a clairvoyant optimum, with the resource being keeps and the constraint being soundness. The factor two with a halved guarantee is the usual shape of such results; the proof here is a few lines because the only structure used is the triangle inequality.

\paragraph{Persistence stability and registration.} The certificate rests on the stability theorem for persistence diagrams \citep{cohensteiner2007}. Registration is Lucas--Kanade \citep{lucas1981,baker2004lk} with a Huber loss, a heuristic that affects only how much is kept, never what is certified. Lipschitz constants of attention \citep{kim2021lipschitz} are what a certificate through a network would need; Section~\ref{sec:shift} says why none is claimed.
